\QHeading{第1题　口算。}
\par\noindent


\begin{oralgrid}
\OralItem{1}{320+250}{570}
\OralItem{2}{65-38}{27}
\OralItem{3}{0+369}{369}
\OralItem{4}{23+17}{40}
\OralItem{5}{35+47}{82}
\OralItem{6}{41+15}{56}
\OralItem{7}{86+49}{135}
\OralItem{8}{72+48}{120}
\OralItem{9}{26+11}{37}
\OralItem{10}{92+77}{169}
\OralItem{11}{27+15}{42}
\OralItem{12}{75+25}{100}\end{oralgrid}
\par

\QHeading{口算思路补充（对应上表题号）}

\Answer{上面列出得数，下面按小题给出不同计算思路。计数单位法适用于小数，整数题可比较凑整、拆分或商不变等方法。}
\MethodTitle{第2小题：$65-38$}

\MethodBlock{口算思路一：补差法}{从减数先补到一个整数，再补到被减数，两段差相加。}{\begin{aligned}[t]&\phantom{={}}65-38\\
&=(40-38)+(65-40)\\
&=2+25\\
&=27\end{aligned}}

\MethodBlock{口算思路二：多减再加回}{}{\begin{aligned}[t]&\phantom{={}}65-38\\
&=65-40+2\\
&=25+2\\
&=27\end{aligned}}
\MethodTitle{第3小题：$0+369$}

\MethodBlock{口算思路一：凑整再减去多加的数}{}{\begin{aligned}[t]&\phantom{={}}0+369\\
&=0+370-1\\
&=370-1\\
&=369\end{aligned}}
\MethodTitle{第7小题：$86+49$}

\MethodBlock{口算思路一：凑整再减去多加的数}{}{\begin{aligned}[t]&\phantom{={}}86+49\\
&=86+50-1\\
&=136-1\\
&=135\end{aligned}}
\MethodTitle{第8小题：$72+48$}

\MethodBlock{口算思路一：凑整再减去多加的数}{}{\begin{aligned}[t]&\phantom{={}}72+48\\
&=72+50-2\\
&=122-2\\
&=120\end{aligned}}


\par\noindent

\QHeading{第2题　填一填。}

\Problem{①\(41+37+13=41+(37+13)\)}
\par

\Problem{②\(x+y=y+x\)}
\Problem{③\(35+(b+65)=(35+65)+b\)}
\Problem{④\(a+b+c=a+c+b\)}
\Problem{⑤\(32+45+55=32+(45+55)\)}
\Problem{⑥\(m+n+t=n+(m+t)\)}
\PlainItem{上面算式用了加法交换律的有（\textcolor{answerred}{2、3、4、6}）。（填序号。）}
\par\noindent

\QHeading{第3题　用简便方法计算下面各题。}

\CalcBlock[方法一：简便计算]{1}{\begin{aligned}[t]&\phantom{={}}\color{black}315+81+219\\&\color{answerred}=315+(81+219)\\&\color{answerred}=315+300\\&\color{answerred}=615\end{aligned}}
\par
\MethodBlock{方法二：先处理括号或乘除，再凑整再减去多加的数}{}{\begin{aligned}[t]&\phantom{={}}315+81+219\\
&=396+219\\
&=396+220-1\\
&=616-1\\
&=615\end{aligned}}


\CalcBlock[方法一：简便计算]{2}{\begin{aligned}[t]&\phantom{={}}\color{black}482+291+118\\&\color{answerred}=(482+118)+291\\&\color{answerred}=600+291\\&\color{answerred}=891\end{aligned}}
\MethodBlock{方法二：先处理括号或乘除，再凑整再减去多加的数}{}{\begin{aligned}[t]&\phantom{={}}482+291+118\\
&=773+118\\
&=773+120-2\\
&=893-2\\
&=891\end{aligned}}

\CalcBlock[方法一：简便计算]{3}{\begin{aligned}[t]&\phantom{={}}\color{black}255+163+237+145\\&\color{answerred}=(255+145)+(163+237)\\&\color{answerred}=400+400\\&\color{answerred}=800\end{aligned}}
\MethodBlock{方法二：交换、结合后先凑整}{保留每个数原来的加减号，先把能凑整的部分放在一起。}{\begin{aligned}[t]&\phantom{={}}255+163+237+145\\
&=(255+145)+163+237\\
&=400+163+237\\
&=563+237\\
&=800\end{aligned}}

\MethodBlock{方法三：按运算顺序直接计算（对照）}{本法只作结果核对；本题要求简便计算，前面的简便解法仍保留。}{\begin{aligned}[t]&\phantom{={}}255+163+237+145\\
&=418+237+145\\
&=655+145\\
&=800\end{aligned}}

\CalcBlock[方法一：简便计算]{4}{\begin{aligned}[t]&\phantom{={}}\color{black}214+172+186\\&\color{answerred}=(214+186)+172\\&\color{answerred}=400+172\\&\color{answerred}=572\end{aligned}}
\MethodBlock{方法二：按运算顺序直接计算（对照）}{本法只作结果核对；本题要求简便计算，前面的简便解法仍保留。}{\begin{aligned}[t]&\phantom{={}}214+172+186\\
&=386+186\\
&=572\end{aligned}}

\CalcBlock[方法一：简便计算]{5}{\begin{aligned}[t]&\phantom{={}}\color{black}154+343+257\\&\color{answerred}=154+(343+257)\\&\color{answerred}=154+600\\&\color{answerred}=754\end{aligned}}
\MethodBlock{方法二：按运算顺序直接计算（对照）}{本法只作结果核对；本题要求简便计算，前面的简便解法仍保留。}{\begin{aligned}[t]&\phantom{={}}154+343+257\\
&=497+257\\
&=754\end{aligned}}

\CalcBlock[方法一：简便计算]{6}{\begin{aligned}[t]&\phantom{={}}\color{black}688+99+112\\&\color{answerred}=(688+112)+99\\&\color{answerred}=800+99\\&\color{answerred}=899\end{aligned}}
\MethodBlock{方法二：按运算顺序直接计算（对照）}{本法只作结果核对；本题要求简便计算，前面的简便解法仍保留。}{\begin{aligned}[t]&\phantom{={}}688+99+112\\
&=787+112\\
&=899\end{aligned}}

\par\noindent

\QHeading{第4题　解决问题。}

\Problem{“唯有牡丹真国色，花开时节动京城。”“五一”假期期间，笑笑一家人去洛阳旅游看牡丹，花费情况如下表，笑笑一家人这次去洛阳旅游共花费多少元？}
\par

\DataTable{4}{吃饭 & 门票 & 购物 & 住宿\\\hline
428元 & 785元 & 572元 & 415元\\\hline}

\MethodTitle{方法一}
\Answer{\Formula{\begin{aligned}[t]&\phantom{={}}428+785+572+415\\&=(428+572)+(785+415)\\&=1000+1200\\&=2200\text{（元）}\end{aligned}}\par 答：共花费2200元。}

\MethodTitle{方法二：凑成整百后统一补偿}
\Answer{428比400多28，785比800少15，572比600少28，415比400多15；多、少正好抵消。}
\Answer{\textbf{多出与不足的数相抵}\par \Formula{\begin{aligned}[t]&\phantom{={}}28+15-15-28\\
&=43-15-28\\
&=28-28\\
&=0\text{（元）}\end{aligned}}}
\Answer{\textbf{按整百数相加}\par \Formula{\begin{aligned}[t]&\phantom{={}}400+800+600+400\\
&=1200+600+400\\
&=1800+400\\
&=2200\text{（元）}\end{aligned}}}
\Answer{答：一共花费2200元。}
